% ============================================================================
% CHAPITRE: TRILEMME CRO CONDITIONNEL CLASSIQUE
% source  : partie2/cro_classique.tex
% Version : 2.2 — alignée sur la formalisation Lean4 (section~\ref{sec:lean4})
% ============================================================================

\section{LE TRILEMME CRO CONDITIONNEL CLASSIQUE}
\label{sec:cro_classique}
\begin{sloppypar}
La section~\ref{sec:clo} a fix\'e le langage, avec trois m\'etriques
d\'efinies via la th\'eorie de l'information~\cite{Shannon48,CoverThomas06} et un Th\'eor\`eme du Goulot d'\'etranglement
Contextuel montrant que $I(V;\Sigma|C) \leq \alpha H(\mathcal{C})$. Aucun r\'esultat n'interdit
pourtant encore d'atteindre simultan\'ement $\Conf = \Rel = \Opp = 0{,}9$. Il
reste donc \`a d\'eterminer o\`u s'arr\^ete le r\'ealisable.

Nous bornons ici la r\'eponse, sous forme du \textbf{Trilemme CRO
conditionnel classique}~\cite{MINKA2025CRO,MINKA2026CRO}. Ce r\'esultat, premi\`ere caract\'erisation quantitative
des limites fondamentales de la cryptographie institutionnelle dans le r\'egime
classique, borne sup\'erieurement le produit $\Conf \cdot \Rel \cdot \Opp$ en fonction de la capacit\'e
d'extraction institutionnelle~$\Delta$, de la granularit\'e des verdicts~$\alpha$ et
du param\`etre de s\'ecurit\'e~$\lambda$. Une forme raffin\'ee de la borne identifie
exactement le r\'egime, d\'etermin\'e par la granularit\'e des verdicts et la richesse
du contexte, dans lequel la contrainte liante tient \`a la capacit\'e d'extraction des interpr\`etes institutionnels, au-del\`a de la seule in\'egalit\'e de Fano.

Deux contributions additionnelles compl\`etent cette section. La classification en
r\'egimes op\'erationnels (sous-section~\ref{subsec:regimes_operationnels}) transforme la
borne abstraite en guide de conception concret selon le niveau d'entropie
contextuelle. Une calibration illustrative (sous-section~\ref{subsec:analyse_comparative})
situe les ordres de grandeur des param\`etres institutionnels \`a partir des trois syst\`emes d\'eploy\'es \'etudi\'es
dans le CLO~\cite{Adida08,Groth16,Johnson01}, qui ne sont pas des instances du th\'eor\`eme
(Remarque~\ref{rem:portee_cas_etude}), conform\'ement \`a la note de m\'ethodologie pr\'esent\'ee dans la sous-section~\ref{subsec:note_methodo}.

Cette section correspond \`a l'\'etape \textbf{B (Borner)} de la progression
F$\to$\textbf{B}$\to$E$\to$S$\to$C$\to$X et apporte la contribution \textbf{C2} de cette th\`ese.

La r\'eponse \`a l'Objectif~\ref{obj:obj2_impossibilite} est donc affirmative, puisque des
limites fondamentales existent dans le r\'egime classique. Leur extension
au contexte quantique sera \'etablie \`a la section suivante.
\end{sloppypar}
\begin{boiteArbitre}
L'arbitre apprend qu'il existe une borne, qui rel\`eve d'une limite fondamentale, ind\'ependante de la technique et repose sur des hypoth\`eses qu'il peut lire et discuter.
Quand ses verdicts sont plus fins que ce qu'il sait extraire, sa propre capacit\'e d'extraction~$\Delta$ le limite davantage que la puissance
de l'adversaire. Les protocoles qu'il examine ne sont donc pas mal con\c{c}us,
car ils se heurtent \`a une incompatibilit\'e structurelle d\'esormais \'enonc\'ee
et prouv\'ee.
\end{boiteArbitre}

% -----------------------------------------------------------------------
\subsection{Mod\`ele formel \'etendu avec param\`etres institutionnels}
\label{subsec:modele_formel_etendu}
% -----------------------------------------------------------------------
\begin{sloppypar}
Nous d\'efinissons trois param\`etres compl\'etant le mod\`ele CLO de la section~\ref{sec:clo}. Ces quantit\'es convertissent des contraintes
qualitatives, \guillemotleft~capacit\'e cognitive limit\'ee~\guillemotright, \guillemotleft~d\'ecisions nuanc\'ees possibles~\guillemotright,
\guillemotleft~protocoles bien con\c{c}us~\guillemotright, en grandeurs mesurables permettant des bornes
rigoureuses~\cite{ClarksonSchneider10,Smith09}.

\subsubsection{Param\`etres institutionnels}
\label{subsubsec:parametres_institutionnels}

\begin{definition}[Capacit\'e d'extraction institutionnelle]\label{def:delta_extraction}

$\Delta \geq 0$ (en bits) quantifie la connaissance pr\'ealable des interpr\`etes
institutionnels et leurs ressources cognitives au-del\`a du contexte formel~$C$~\cite{Simon82,Miller56}.
Ce param\`etre agr\`ege les pr\'ec\'edents l\'egaux et la jurisprudence applicable~\cite{Mason07,Kerr04},
l'expertise de domaine des d\'ecideurs~\cite{Ericsson95,Chi2006}, la documentation et les proc\'edures
\'etablies, les capacit\'es de d\'elib\'eration collaborative entre experts~\cite{Sunstein2005,Hong2016}, ainsi
que la m\'emoire institutionnelle et les bases de donn\'ees historiques~\cite{Argyris1978,Walsh1995}.

\end{definition}

\begin{definition}[Granularit\'e des verdicts]\label{def:alpha_granularite}

$\alpha > 0$ (sans dimension) mesure la finesse de distinction des verdicts
institutionnels relativement \`a la richesse contextuelle~:
\begin{equation}
\log_2 |\mathcal{V}_\lambda| \leq \alpha \cdot H(C)
\label{eq:alpha_definition}
\end{equation}
Un tribunal rendu apte \`a distinguer douze niveaux de fraude dans un contexte \`a
entropie $H(C) = 4$ bits correspond \`a $\alpha = \log_2 12/4 \approx 0{,}9$~\cite{Shannon48,CoverThomas06}.

\end{definition}

\begin{definition}[R\'egularit\'e du protocole]\label{def:beta_regularite}

$\beta \geq 1$ borne la divergence Kullback-Leibler entre distributions de
preuves pour messages distincts~\cite{Kullback1951,CoverThomas06}~:
\begin{equation}
\KL\!\left(P_{\Sigma|m_0,c} \,\|\, P_{\Sigma|m_1,c}\right) \leq \beta \cdot I(M;\Sigma|C)
\label{eq:beta_definition}
\end{equation}
Ce param\`etre \'ecarte les protocoles pathologiques dont la divergence KL exploserait
pour certaines paires de messages~\cite{Tsybakov2009}. Au sens statistique, un protocole
qui publie le message en clair (signature) a des lois de preuve \`a supports disjoints
et $\KL = \infty$, de sorte qu'il n'est pas $\beta$-r\'egulier. La $\beta$-r\'egularit\'e est une
hypoth\`ese sur la classe de protocoles vis\'ee (protocoles \`a fuite contr\^ol\'ee), non
une propri\'et\'e observ\'ee des protocoles d\'eploy\'es~\cite{Pedersen91,Groth16}.

\end{definition}
\end{sloppypar}

% -----------------------------------------------------------------------
\subsubsection{Hypoth\`eses sur les capacit\'es institutionnelles}
\label{subsubsec:hypotheses_capacite}
% -----------------------------------------------------------------------

\begin{sloppypar}
Nous posons quatre hypoth\`eses sur les capacit\'es
institutionnelles~\cite{Simon1955,Miller56,Ericsson95}. Chacune repose
sur des justifications multidisciplinaires ; aucune n'a \'et\'e valid\'ee
\`a ce stade par protocoles exp\'erimentaux contr\^ol\'es mesurant
directement $\Delta$, $\alpha$, $\beta$ aupr\`es d'acteurs
institutionnels r\'eels, cette limitation est reconnue explicitement
(sous-section~\ref{subsec:note_methodo}).

\begin{definition}[Interpr\'etation institutionnelle admissible]\label{def:interp_admissible}
Une interpr\'etation est admissible si le verdict s'\'ecrit
$V = \mathsf{Interpret}(\Sigma, C, K)$, o\`u $K$ repr\'esente les ressources
propres de l'institution (jurisprudence, expertise, d\'elib\'eration,
al\'eas de d\'ecision) et satisfait
\begin{equation}
K \perp M \mid (\Sigma, C).
\label{eq:interp_admissible}
\end{equation}
L'institution n'acc\`ede au message que par la preuve, de sorte que $M \to (\Sigma,C) \to V$
est une cha\^ine de Markov. Le verdict peut \^etre stochastique.
\end{definition}

\begin{assumption}[Extraction born\'ee]
\label{ass:extraction_bornee}
Pour tout protocole CLO :
\begin{equation}
I(V;\Sigma|C) \leq I(M;\Sigma|C) + \Delta
\label{eq:extraction_bornee}
\end{equation}
\end{assumption}

\begin{proposition}[Fondement informationnel de l'Hypoth\`ese~\ref{ass:extraction_bornee}]
\label{prop:fondement_extraction}
Pour toute interpr\'etation admissible,
\begin{equation}
I(V;\Sigma|C) = I(V;M|C) + I(V;\Sigma|M,C) \;\leq\; I(M;\Sigma|C) + I(V;\Sigma|M,C).
\label{eq:fondement_extraction}
\end{equation}
En cons\'equence, l'Hypoth\`ese~\ref{ass:extraction_bornee} est satisfaite
d\`es que
\begin{equation}
I(V;\Sigma|M,C) \leq \Delta,
\label{eq:extraction_bornee_forte}
\end{equation}
et $\Delta$ s'interpr\`ete exactement comme l'information que
l'institution tire de la preuve \emph{au-del\`a du message}.
\end{proposition}

\begin{Preuve}[Proposition~\ref{prop:fondement_extraction}]\label{prv:fondement_extraction}
\begin{cmtEtape}{Objectif.}
Justifier l'Hypoth\`ese~\ref{ass:extraction_bornee} \`a partir d'une
propri\'et\'e structurelle du processus d'interpr\'etation.
\end{cmtEtape}
Par la r\`egle de cha\^ine appliqu\'ee \`a $I(V;M,\Sigma|C)$ de deux
mani\`eres, on obtient
\begin{equation}
I(V;M|C) + I(V;\Sigma|M,C) = I(V;\Sigma|C) + I(V;M|\Sigma,C).
\end{equation}
\cmtFormule{Deux d\'ecompositions de la m\^eme information mutuelle.}
Sous~\eqref{eq:interp_admissible}, $V$ ne d\'epend de $M$ qu'\`a travers
$(\Sigma,C)$, donc $I(V;M|\Sigma,C)=0$, ce qui donne l'\'egalit\'e
de~\eqref{eq:fondement_extraction}. L'in\'egalit\'e de traitement des
donn\'ees sur $M \to (\Sigma,C) \to V$ donne
$I(V;M|C) \leq I(M;\Sigma|C)$.
\cmt{L'institution ne peut pas apprendre sur le message plus que la preuve n'en contient~; elle peut en revanche lire dans la preuve des traits qui ne sont pas le message.}
\end{Preuve}

\begin{assumption}[Verdict efficace]
\label{ass:verdicts_efficace}
La taille de l'espace de verdict satisfait $|\mathcal{V}_\lambda| \leq
2^{\alpha H(C)}$ pour une constante $\alpha > 0$, \emph{et} la fonction
de verdict utilise efficacement le contexte disponible :
\begin{equation}
H(V|C) \geq \alpha H(C) - 1 =: \gamma, \qquad \gamma \geq 1
\label{eq:verdicts_efficace}
\end{equation}
\end{assumption}

\begin{remark}[Statut de l'Hypoth\`ese~\ref{ass:verdicts_efficace}]
\label{rem:statut_verdict_efficace}
L'in\'egalit\'e $|\mathcal{V}_\lambda| \leq 2^{\alpha H(C)}$ est une
borne \emph{sup\'erieure} sur la taille de l'espace de verdict. Elle ne
permet pas, \`a elle seule, de d\'eduire la borne \emph{inf\'erieure}
$H(V|C) \geq \alpha H(C) - 1$. Cette borne inf\'erieure est la
propri\'et\'e d'\textbf{efficacit\'e} du verdict, qui doit \^etre
postul\'ee explicitement. Elle formalise l'id\'ee qu'une institution
exploitant effectivement le contexte disponible distingue au moins
$\gamma$ bits de verdicts. Pour les param\`etres institutionnels typiques
($H(C) \geq 10$ bits, $\alpha \approx 1$), on a $\gamma \geq 9$.
\end{remark}

\begin{assumption}[R\'egularit\'e]
\label{ass:regularite}
La divergence KL satisfait $\KL(P_{\Sigma|m_0,c} \| P_{\Sigma|m_1,c})
\leq \beta \cdot I(M;\Sigma|C)$.
\end{assumption}

\begin{assumption}[R\'egime de pr\'eservation de la confidentialit\'e]
\label{ass:regime}
Il existe une constante $\kappa \geq 0$, ind\'ependante de $\lambda$,
telle que :
\begin{equation}
I(M;\Sigma|C) \leq \frac{\kappa}{\lambda}
\label{eq:regime}
\end{equation}
\end{assumption}

\begin{remark}[Statut et r\^ole de l'Hypoth\`ese~\ref{ass:regime}]
\label{rem:statut_regime}
Dans la version initiale de la preuve, ce r\'egime \'etait pr\'esent\'e
comme une \emph{cons\'equence} de l'optimisation de $f(I)$ sur
$[0, \lambda]$. La formalisation a conduit \`a une r\'esolution alternative de cette
optimisation. La d\'eriv\'ee $f'(x) = 2x - \frac{5}{2}ax^{3/2}$ s'annule en
$\sqrt{x^*} = \frac{4}{5a}$ et ce point est un \emph{maximiseur} de $f$. Le r\'egime
$I = O(1/\lambda)$ ne peut donc pas \^etre d\'eduit de cette
optimisation.

L'Hypoth\`ese~\ref{ass:regime} est donc pos\'ee explicitement
comme condition du th\'eor\`eme. Elle est pr\'ecis\'ement le r\'egime
dans lequel les protocoles $\beta$-r\'eguliers garantissent la
confidentialit\'e, comme l'\'etablit le
Lemme~\ref{lem:confidentiality_guarantee}. Si
$I \leq \kappa/\lambda$, alors
$\Conf \geq 1 - \sqrt{c_0 \beta \kappa/\lambda} \to 1$ quand
$\lambda \to \infty$.
\end{remark}
\end{sloppypar}

% -----------------------------------------------------------------------
\begin{figure}[htbp]
\centering
\begin{tikzpicture}[scale=1.0, transform shape]
% --- Borne de Fano (1+kappa/lambda)/lambda ---
\draw[enspred, thick, dashed]
    plot[smooth] coordinates {
        (1.60, 6.348)
        (2.00, 5.062)
        (2.40, 4.210)
        (2.80, 3.603)
        (3.20, 3.149)
        (4.00, 2.516)
        (4.80, 2.094)
        (5.60, 1.794)
        (6.40, 1.569)
        (8.00, 1.254)
        (9.60, 1.044)
    }
    node[right, font=\small, text=enspred] {Fano seul};
% --- Borne raffin\'ee, Delta < gamma ---
\draw[enspblue, thick]
    plot[smooth] coordinates {
        (1.60, 1.676)
        (2.00, 1.336)
        (2.40, 1.110)
        (2.80, 0.950)
        (3.20, 0.830)
        (4.00, 0.663)
        (4.80, 0.552)
        (5.60, 0.472)
        (6.40, 0.413)
        (8.00, 0.330)
        (9.60, 0.275)
    }
    node[right, font=\small, text=enspblue, align=left] {Borne raffin\'ee\\$\Delta=5$, $\gamma=19$};
% --- Axes ---
\draw[-{Stealth[length=3mm]}, thick] (0,0) -- (10.5, 0)
    node[right, font=\small] {$\lambda$};
\draw[-{Stealth[length=3mm]}, thick] (0,0) -- (0, 7.0)
    node[above, font=\small] {Majorant de $\Conf \cdot \Rel \cdot \Opp$};
\foreach \y/\lab in {2.0/0{,}005, 4.0/0{,}010, 6.0/0{,}015} {
    \draw[thin, chaptercolor] (-0.1, \y) -- (0.1, \y);
    \node[left, font=\footnotesize] at (-0.1, \y) {$\lab$};
}
\node[left, font=\footnotesize] at (-0.1, 0) {$0$};
\foreach \x/\lab in {1.6/64, 3.2/128, 6.4/256, 9.6/384} {
    \draw[thin, chaptercolor] (\x, -0.1) -- (\x, 0.1);
    \node[below, font=\footnotesize] at (\x, -0.1) {$\lab$};
}
\end{tikzpicture}
\caption[Borne raffin\'ee du Trilemme CRO et borne de Fano]{%
  Majorants de $\Conf \cdot \Rel \cdot \Opp$ en fonction de $\lambda$
  ($\kappa = 1$)~: borne raffin\'ee~\eqref{eq:trilemme_cro_raffine} pour
  $\Delta = 5$, $\alpha = 1$, $H(\mathcal{C}) = 20$ bits ($\gamma = 19$, trait
  plein) et borne de Fano $(1+\kappa/\lambda)/\lambda$ (tirets). L'\'ecart entre
  les deux mesure l'effet de la capacit\'e institutionnelle lorsque
  $\Delta < \gamma$~; pour $\Delta \geq \gamma$, les deux courbes co\"incident.%
}\label{fig:cro_bound}
\end{figure}

% -----------------------------------------------------------------------
\subsection{\'Enonc\'e et preuve du Trilemme CRO classique}
\label{subsec:enonce_trilemme}
% -----------------------------------------------------------------------
\begin{sloppypar}
Nous \'etablissons le r\'esultat principal du r\'egime classique.

\begin{remark}[Version de la preuve]
\label{rem:version_preuve_classique}
La version de la preuve pr\'esent\'ee ici diff\`ere de celle de l'article
publi\'e~\cite{MINKA2026CRO}. Les formulations alternatives retenues, issues
de la certification Lean4, sont d\'etaill\'ees \`a la
section~\ref{subsec:lean4_erreurs}.
\end{remark}

\begin{lemma}[Garantie de confidentialit\'e pour les protocoles r\'eguliers]
\label{lem:confidentiality_guarantee}
Pour tout protocole $\beta$-r\'egulier (Hypoth\`ese~\ref{ass:regularite}) :
\begin{equation}
\Conf(\mathcal{P}) \geq 1 - \sqrt{c_0 \beta \cdot I(M;\Sigma|C)}
\label{eq:lemma4}
\end{equation}
o\`u $c_0 = 1/(2\ln 2) \approx 0{,}721$ est la constante de Pinsker.
\end{lemma}

\begin{proof}
Par le Lemme~\ref{lem:compromis_conf_mi} (Pinsker) et
l'Hypoth\`ese~\ref{ass:regularite}, pour tous $m_0, m_1 \in
\mathcal{M}_\lambda$ et $c \in \mathcal{C}_\lambda$, on a
\begin{equation}
\TV(P_{\Sigma|m_0,c}, P_{\Sigma|m_1,c}) \leq \sqrt{c_0 \KL(P_{\Sigma|m_0,c} \| P_{\Sigma|m_1,c})} \leq \sqrt{c_0 \beta \cdot I(M;\Sigma|C)}
\end{equation}
En prenant le maximum sur $(m_0, m_1, c)$ et en appliquant la
D\'efinition~\ref{def:confidentialite_clo}, on obtient~\eqref{eq:lemma4}.
\end{proof}

\begin{remark}[Interpr\'etation du Lemme~\ref{lem:confidentiality_guarantee}]
\label{rem:interp_lemme4}
Le Lemme~\ref{lem:confidentiality_guarantee} \'enonce que la faible
information mutuelle conditionnelle, combin\'ee \`a la
$\beta$-r\'egularit\'e, \emph{garantit} la confidentialit\'e. Si
$I(M;\Sigma|C) \leq \kappa/\lambda$, alors
$\Conf(\mathcal{P}) \geq 1 - \sqrt{c_0 \beta \kappa/\lambda} \to 1$
quand $\lambda \to \infty$.
\end{remark}

\begin{remark}[Constante $c_0$]
\label{rem:constante_c0}
Avec des divergences en bits, la constante de Pinsker serr\'ee est
$(\ln 2)/2 \approx 0{,}347$. La constante $c_0 = 1/(2\ln 2) \approx
0{,}721$ utilis\'ee dans ce chapitre, reprise de~\cite{MINKA2026CRO}, est plus
grande. Toute in\'egalit\'e \'enonc\'ee avec $c_0$ reste donc valide sans \^etre
serr\'ee. Le qualificatif \og serr\'ee\fg{} employ\'e dans l'article publi\'e
s'entend en ce sens.
\end{remark}

\begin{theorem}[Trilemme CRO conditionnel classique]
\label{thm:trilemme_cro_conditionnel}
Sous les Hypoth\`eses~\ref{ass:extraction_bornee},~\ref{ass:verdicts_efficace},
\ref{ass:regularite} et~\ref{ass:regime}, pour tout protocole CLO avec
$\lambda \geq 64$, $H(M) = \lambda$, $M \perp C$~:
\begin{equation}
\Conf \cdot \Rel \cdot \Opp \leq \frac{\Delta}{\lambda} + O(\lambda^{-2}) + \negl(\lambda)
\label{eq:trilemme_cro_conditionnel}
\end{equation}
Plus pr\'ecis\'ement, avec $\gamma = \alpha H(\mathcal{C}) - 1$~:
\begin{equation}
\Conf \cdot \Rel \cdot \Opp \leq \frac{1}{\gamma}\left(\frac{\Delta}{\lambda}
+ \frac{\kappa(\Delta+1)}{\lambda^2} + \frac{\kappa^2}{\lambda^3}\right)
\label{eq:trilemme_cro_explicite}
\end{equation}
et, sous la forme la plus pr\'ecise,
\begin{equation}
\Conf \cdot \Rel \cdot \Opp \leq \frac{1+\kappa/\lambda}{\lambda}\,
  \min\left\{1,\; \frac{\kappa/\lambda + \Delta}{\gamma}\right\}.
\label{eq:trilemme_cro_raffine}
\end{equation}
De plus, $\Conf(\mathcal{P}) \geq 1 - \sqrt{c_0 \beta \kappa/\lambda}$
par le Lemme~\ref{lem:confidentiality_guarantee}.
\end{theorem}

La preuve proc\`ede en cinq \'etapes. Les deux premi\`eres construisent les bornes
individuelles sur $\Rel$ et $\Conf$ via les in\'egalit\'es classiques de Fano et de
Pinsker~\cite{Fano61,CoverThomas06,Pinsker1964,Tsybakov2009}~; les \'etapes trois et quatre bornent $\Opp$ par les hypoth\`eses
institutionnelles~\cite{Simon82,Miller56}~; la derni\`ere combine le produit et applique
le r\'egime institutionnel.

\begin{Preuve}[Preuve d\'etaill\'ee]\label{prv:cro_classique}

\begin{cmtEtape}{Vue d'ensemble de la preuve.}
Posons $x = I(M;\Sigma|C)$, l'information que la preuve r\'ev\`ele sur le
message une fois le contexte connu. Toute la d\'emonstration consiste \`a
exprimer les trois m\'etriques en fonction de cette \emph{unique} quantit\'e,
puis \`a constater qu'elles ne peuvent cro\^itre ensemble. La fiabilit\'e
exige $x$ grand~; la confidentialit\'e exige $x$ petit~; l'opposabilit\'e,
elle, est plafonn\'ee par ce que l'institution parvient \`a extraire. Elle se
d\'eroule en cinq \'etapes, o\`u (1)~Fano borne $\Rel$ par $x$~; (2)~Pinsker et la r\'egularit\'e
garantissent $\Conf$ \`a partir de $x$~; (3)~l'extraction born\'ee borne
$\Opp$~; (4)~la borne sur les verdicts ach\`eve cette contrainte~;
(5)~on forme le produit et on applique le r\'egime.
\end{cmtEtape}

    \paragraph{\'Etape 1. Fano contraint la fiabilit\'e.}
    L'outil mobilis\'e ici est l'in\'egalit\'e de Fano, introduite \`a la
    sous-sous-section~\ref{subsubsec:fano_classique}. Elle \'enonce que, si l'on
    cherche \`a deviner une variable al\'eatoire \`a partir d'une observation,
    l'incertitude qui subsiste sur cette variable est born\'ee par la
    probabilit\'e de se tromper. Appliquons-la au message.

    Soit $\hat{M}$ l'estimateur optimal de $M$ construit \`a partir du couple
    $(\Sigma, C)$, c'est-\`a-dire la meilleure reconstruction possible du
    message \`a partir de la preuve et du contexte. Notons
    $p_e = 1 - \Rel$ sa probabilit\'e d'erreur. Fano donne alors
    \begin{equation}
    H(M|\Sigma,C) \leq h(p_e) + p_e\log_2(2^\lambda - 1)
    \label{eq:fano_applique}
    \end{equation}
    \cmtFormule{L'incertitude r\'esiduelle sur $M$ est born\'ee par la probabilit\'e d'erreur.}
    \cmt{$h(\cdot)$ d\'esigne l'entropie binaire, $h(p) = -p\log_2 p - (1-p)\log_2(1-p)$.}
    \cmt{Le second terme compte le nombre de messages restants \`a d\'epartager en cas d'erreur, soit $2^\lambda - 1$ messages.}

    L'entropie binaire ne d\'epasse jamais un bit, $h(p_e) \leq 1$, tandis que
    $\log_2(2^\lambda - 1) \leq \lambda$. En rempla\c{c}ant dans
    \eqref{eq:fano_applique}, on obtient
    \begin{equation}
    H(M|\Sigma,C) \leq 1 + (1 - \Rel)\,\lambda
    \end{equation}
    \cmt{Aucune restriction du type $\Rel \geq 1/2$ n'est n\'ecessaire ici, car la majoration $h(p_e) \leq 1$ vaut pour toute valeur de $p_e$.}

    Les messages \'etant uniformes sur $\mathcal{M}_\lambda$, leur entropie
    vaut $H(M) = \lambda$ (sous-sous-section~\ref{subsubsec:entropie_info_mutuelle}).
    L'information mutuelle \'etant la diff\'erence entre l'entropie et
    l'entropie conditionnelle, il vient
    \begin{equation}
    I(M;\Sigma,C) = H(M) - H(M|\Sigma,C) \geq \lambda - 1 - (1-\Rel)\lambda = \Rel \cdot \lambda - 1
    \label{eq:fano_mi}
    \end{equation}
    \cmtFormule{L'information extraite sur $M$ vaut au moins $\Rel\lambda - 1$.}

    Le protocole est con\c{c}u de sorte que le message soit ind\'ependant du
    contexte institutionnel, $M \perp C$~; autrement dit, savoir dans quelle juridiction on
    se trouve n'apprend rien sur le message transmis. Par cons\'equent
    $I(M;\Sigma,C) = I(M;\Sigma|C) = x$, d'o\`u
    \begin{equation}
    \boxed{\;\Rel \leq \frac{x + 1}{\lambda}\;}
    \label{eq:etape1_rel}
    \end{equation}
    \cmt{Lecture. Une preuve qui ne r\'ev\`ele presque rien sur le message ($x$ petit) ne permet pas de le reconstituer, donc la fiabilit\'e chute.}

    \paragraph{\'Etape 2. Pinsker et r\'egularit\'e garantissent la confidentialit\'e.}
    On mobilise cette fois l'in\'egalit\'e de Pinsker
    (sous-sous-section~\ref{subsubsec:pinsker_classique}), qui relie deux
    fa\c{c}ons de mesurer l'\'ecart entre deux lois. La premi\`ere est la distance
    en variation totale, qui gouverne la capacit\'e d'un adversaire \`a les
    distinguer~; la seconde est la divergence de Kullback-Leibler
    (sous-sous-section~\ref{subsubsec:distances_divergences}).

    Pour tout couple de messages $(m_0, m_1)$ et tout contexte $c$, Pinsker
    donne $\TV^2 \leq c_0 \KL$ avec $c_0 = 1/(2\ln 2) \approx 0{,}721$, tandis que
    l'Hypoth\`ese~\ref{ass:regularite} de $\beta$-r\'egularit\'e majore cette
    divergence par $\beta\,x$. En composant les deux, puis en passant au
    maximum sur les couples et les contextes, on obtient le
    Lemme~\ref{lem:confidentiality_guarantee}, soit
    \begin{equation}
    \boxed{\;\Conf \geq 1 - \sqrt{c_0 \beta x}\;}
    \label{eq:etape2_conf}
    \end{equation}
    \cmtFormule{Confidentialit\'e minor\'ee par une fonction d\'ecroissante de $x$.}
    \cmt{Il s'agit d'une \emph{garantie}, puisque tant que la preuve divulgue peu ($x \leq \kappa/\lambda$), la confidentialit\'e tend vers $1$.}
    \cmt{Cette \'etape ne sert pas \`a majorer le produit, mais \`a montrer que le r\'egime consid\'er\'e est bien celui o\`u la confidentialit\'e est pr\'eserv\'ee.}

    \paragraph{\'Etape 3. L'extraction born\'ee contraint l'opposabilit\'e.}
    Par d\'efinition, l'opposabilit\'e est la fraction de l'incertitude
    institutionnelle que la preuve r\'esout,
    $\Opp = I(V;\Sigma|C)/H(V|C)$. Son num\'erateur mesure ce que la preuve
    apporte au verdict, son d\'enominateur l'ampleur de la d\'ecision \`a
    prendre.

    L'Hypoth\`ese~\ref{ass:extraction_bornee} borne ce num\'erateur,
    puisque l'institution ne tire de la preuve que ce que celle-ci contient sur le
    message, $x$, augment\'e d'au plus $\Delta$ bits qu'elle extrait par ses
    propres moyens (pr\'ec\'edents, expertise, proc\'edures). Il vient
    \begin{equation}
    \Opp = \frac{I(V;\Sigma|C)}{H(V|C)} \leq \frac{x + \Delta}{H(V|C)}
    \label{eq:etape3_opp}
    \end{equation}
    \cmtFormule{Le num\'erateur est major\'e par $x + \Delta$.}
    \cmt{$\Delta$ est une capacit\'e cognitive, born\'ee par la rationalit\'e de l'interpr\`ete~; elle ne cro\^it pas avec le param\`etre de s\'ecurit\'e.}

    \paragraph{\'Etape 4. La borne sur les verdicts compl\`ete la contrainte sur $\Opp$.}
    Il reste \`a minorer le d\'enominateur. L'Hypoth\`ese~\ref{ass:verdicts_efficace}
    \'enonce qu'un verdict institutionnel exploite effectivement le contexte
    dont il dispose, ce qui s'\'ecrit
    \begin{equation}
    H(V|C) \geq \gamma := \alpha H(\mathcal{C}) - 1 \geq 1
    \label{eq:etape4_hvc}
    \end{equation}
    \cmtFormule{L'incertitude sur les verdicts est minor\'ee par $\gamma$.}
    \cmt{Un verdict d\'eg\'en\'er\'e, toujours identique, aurait $H(V|C) = 0$ et rendrait $\Opp$ ind\'efini, cas que l'hypoth\`ese \'ecarte.}

    En injectant \eqref{eq:etape4_hvc} dans \eqref{eq:etape3_opp}, on obtient
    \begin{equation}
    \boxed{\;\Opp \leq \frac{x + \Delta}{\gamma}\;}
    \label{eq:etape4_opp}
    \end{equation}

    \paragraph{\'Etape 5. Formation du produit et application du r\'egime.}
    Les trois m\'etriques sont d\'esormais exprim\'ees en fonction de $x$. En
    multipliant \eqref{eq:etape1_rel} et \eqref{eq:etape4_opp} puis en majorant
    trivialement $\Conf \leq 1$, on obtient
    \begin{equation}
    \Conf \cdot \Rel \cdot \Opp \leq \frac{x+1}{\lambda} \cdot \frac{x+\Delta}{\gamma}
    \label{eq:produit_brut}
    \end{equation}
    \cmt{On majore ici $\Conf$ par $1$ sans recourir \`a \eqref{eq:etape2_conf}, car dans le r\'egime consid\'er\'e la confidentialit\'e est proche de $1$ et la contrainte vient des deux autres facteurs.}

    Le membre de droite, produit de deux fonctions croissantes et
    positives de $x$, est donc croissant en $x$. Il suffit ainsi de
    l'\'evaluer \`a la borne sup\'erieure du r\'egime. Par
    l'Hypoth\`ese~\ref{ass:regime}, $x \leq \kappa/\lambda$, d'o\`u
    \begin{equation}
    \Conf \cdot \Rel \cdot \Opp \leq \frac{1}{\gamma}\left(\frac{\kappa/\lambda + 1}{\lambda}\right)\left(\frac{\kappa}{\lambda} + \Delta\right)
    \end{equation}
    \cmt{Substitution directe de $x = \kappa/\lambda$ dans \eqref{eq:produit_brut}.}

    Le d\'eveloppement du produit donne quatre termes, que l'on ordonne par
    puissance d\'ecroissante de $\lambda$~:
    \begin{align}
    \left(\frac{\kappa/\lambda + 1}{\lambda}\right)\left(\frac{\kappa}{\lambda} + \Delta\right)
    &= \frac{\kappa^2}{\lambda^3} + \frac{\kappa\Delta}{\lambda^2}
     + \frac{\kappa}{\lambda^2} + \frac{\Delta}{\lambda} \nonumber \\
    &= \frac{\Delta}{\lambda} + \frac{\kappa(\Delta + 1)}{\lambda^2} + \frac{\kappa^2}{\lambda^3}
    \end{align}
    \cmtFormule{Les deux termes crois\'es en $\kappa$ se regroupent en $\kappa(\Delta+1)/\lambda^2$.}

    On obtient ainsi la forme explicite annonc\'ee
    \begin{equation}
    \Conf \cdot \Rel \cdot \Opp \leq \frac{1}{\gamma}\left(\frac{\Delta}{\lambda} + \frac{\kappa(\Delta+1)}{\lambda^2} + \frac{\kappa^2}{\lambda^3}\right)
    \end{equation}

    Comme $\gamma \geq 1$, diviser par $\gamma$ ne fait que renforcer la
    majoration. Les deux derniers termes \'etant en $O(\lambda^{-2})$, il vient
    \begin{equation}
    \Conf \cdot \Rel \cdot \Opp \leq \frac{\Delta}{\lambda} + O(\lambda^{-2}) + \negl(\lambda)
    \end{equation}
    \cmtFormule{Le terme dominant est $\Delta/\lambda$.}
    \cmt{Lecture finale. Tant que $\Delta < \gamma$, la contrainte qui plafonne le produit est la capacit\'e institutionnelle~; au-del\`a, c'est la borne de Fano (Remarque~\ref{rem:activite_delta}).}

    \textbf{Forme raffin\'ee.} Comme $\Opp \leq 1$
    (Th\'eor\`eme~\ref{thm:bornes_mesures}), l'\'Etape~2 et l'Hypoth\`ese~\ref{ass:regime}
    donnent aussi $\Opp \leq \min\{1, (\kappa/\lambda+\Delta)/\gamma\}$~;
    combin\'ee \`a l'\'Etape~1 sous l'Hypoth\`ese~\ref{ass:regime},
    ($\Rel \leq (1+\kappa/\lambda)/\lambda$) ainsi qu'\`a $\Conf \leq 1$, elle
    donne~\eqref{eq:trilemme_cro_raffine}. Comme le minimum est major\'e par
    son second argument, la forme raffin\'ee implique la forme explicite.

\end{Preuve}

\begin{remark}[Quand la capacit\'e institutionnelle est-elle contraignante~?]
\label{rem:activite_delta}
Le facteur $\min\{1,(\kappa/\lambda+\Delta)/\gamma\}$ est strictement
inf\'erieur \`a~$1$ si et seulement si
\begin{equation}
\Delta < \gamma - \frac{\kappa}{\lambda}
\;\Longleftrightarrow\;
H(\mathcal{C}) > H^\sharp := \frac{\Delta + 1 + \kappa/\lambda}{\alpha}.
\label{eq:seuil_activite}
\end{equation}
Dans ce r\'egime (\emph{contexte riche relativement \`a la capacit\'e}), la
borne est port\'ee par le rapport $\Delta/\gamma$ et la capacit\'e
institutionnelle est la contrainte liante. Dans le r\'egime compl\'ementaire
($H(\mathcal{C}) \leq H^\sharp$), le produit est born\'e par la seule
in\'egalit\'e de Fano, $\Relrec \leq (1+\kappa/\lambda)/\lambda$~; la
forme~\eqref{eq:trilemme_cro_explicite} n'apporte alors pas d'information
suppl\'ementaire. Si le verdict \'etait une fonction d\'eterministe de
$(\Sigma,C)$, les Hypoth\`eses~\ref{ass:extraction_bornee}
et~\ref{ass:verdicts_efficace} imposeraient $\Delta \geq \gamma - I(M;\Sigma|C)$,
de sorte qu'on serait toujours dans le second r\'egime
(Remarque~\ref{rem:verdict_deterministe}). Ces faits sont certifi\'es
en Lean (\texttt{trilemma\_refined}, \texttt{refined\_bound\_is\_fano\_when\_delta\_ge\_gamma},
\texttt{deterministic\_forces\_capacity}, section~\ref{sec:lean4}).
\end{remark}

\begin{remark}[Port\'ee du Th\'eor\`eme~\ref{thm:trilemme_cro_conditionnel}]
\label{rem:scope_trilemme_classique}
L'Hypoth\`ese~\ref{ass:regime} est n\'ecessaire. Sans borne sur
$I(M;\Sigma|C)$, le produit $\Conf \cdot \Rel \cdot \Opp$ n'a pas besoin
de d\'ecro\^itre avec $\lambda$. En effet, le protocole qui r\'ev\`ele $M$ avec
probabilit\'e $1/2$ et produit sinon une preuve uniform\'ement
al\'eatoire est $\beta$-r\'egulier avec $\beta \leq 1{,}04$ pour tout
$\lambda \geq 64$, satisfait les
Hypoth\`eses~\ref{ass:extraction_bornee}--\ref{ass:regularite} et a
$\Conf = 1/2$, $\Rel = 1/2 + 2^{-\lambda-1}$ et, avec un contexte $C$
form\'e de deux bits uniformes ind\'ependants de $(M,\Sigma)$, $\alpha = 1$
(donc $\gamma = \alpha H(C) - 1 = 1$) et $V$ le premier bit de $\Sigma$
(donc $H(V|C) = 1 \geq \gamma$ et $\log_2|\mathcal{V}_\lambda| = 1 \leq \alpha H(C)$),
$\Opp = 1$, pour tout $\lambda$. L'Hypoth\`ese~\ref{ass:extraction_bornee} est
satisfaite avec $\Delta = 0$, puisque $I(V;\Sigma|C) = 1 \leq I(M;\Sigma|C)$. Ici
$I(M;\Sigma|C) = \lambda/2 - 1 + O(\lambda 2^{-\lambda})$ se situe hors
du r\'egime~\eqref{eq:regime}.
\end{remark}

\begin{remark}[Statut de v\'erification]
\label{rem:statut_verification_classique}
Le noyau alg\'ebrique de cette preuve est certifi\'e m\'ecaniquement en
Lean4 \`a la section~\ref{sec:lean4}. Les in\'egalit\'es de Fano et de
Pinsker, ainsi que les hypoth\`eses institutionnelles, y sont
formalis\'ees comme propri\'et\'es des instances de protocole. Le
d\'etail des formulations alternatives issues de la certification est pr\'esent\'e \`a la
section~\ref{subsec:lean4_erreurs}.
\end{remark}

\begin{corollary}[Forme simplifi\'ee]\label{cor:forme_simplifiee}

Pour $\lambda \geq 128$, $\beta = O(1)$, $\gamma = O(1)$~\cite{Goldreich01,KatzLindell14}~:
\begin{equation}
\Conf \cdot \Rel \cdot \Opp \lesssim \frac{\Delta}{\lambda}
\label{eq:trilemme_simplifie}
\end{equation}

\end{corollary}

Pour $\Delta = 5$ bits, $\alpha = 1$, $H(\mathcal{C}) = 20$ bits
($\gamma = 19$), $\kappa = 1$ et $\lambda = 128$, la
forme~\eqref{eq:trilemme_cro_raffine} donne
$\frac{1+1/128}{128}\cdot\frac{1/128 + 5}{19} \approx 2{,}1\cdot10^{-3}$,
contre $\approx 7{,}9\cdot10^{-3}$ pour la borne de Fano seule, de sorte que la capacit\'e
institutionnelle divise la borne par pr\`es de quatre. Pour $\Delta = 35$
bits dans le m\^eme contexte, $\Delta > \gamma$ et la borne retombe sur
celle de Fano ($\approx 7{,}9\cdot10^{-3}$)~; la forme publi\'ee
$\Delta/\lambda \approx 0{,}27$ est alors tr\`es l\^ache. Ces valeurs ont
statut d'illustration num\'erique~; elles sont calcul\'ees sous les hypoth\`eses
institutionnelles et ne constituent pas des mesures empiriques~\cite{Tsybakov2009}.

\textbf{Augmenter $\lambda$ n'am\'eliore pas le produit CRO si $\Delta$ reste fixe.}
\`A hypoth\`eses fix\'ees, la borne d\'ecro\^it en $1/\lambda$ parce que
l'Hypoth\`ese~\ref{ass:regime} force la fiabilit\'e de reconstruction \`a
d\'ecro\^itre~; la capacit\'e institutionnelle d\'etermine, dans le
r\'egime~\eqref{eq:seuil_activite}, la fraction de ce plafond que
l'institution peut exploiter. C'est l\`a que les limites de l'interpr\`ete, au-del\`a
de la seule puissance des adversaires, deviennent la contrainte liante.

\begin{boiteArbitre}
L'arbitre comprend que commander un algorithme plus fort ne changera rien,
puisque le goulot d'\'etranglement se situe en lui d\`es que ses verdicts
sont plus fins que ce qu'il sait extraire.
\end{boiteArbitre}

\begin{figure}[htbp]
\centering
\begin{tikzpicture}[
    scale=0.8,
    transform shape,
    every node/.style={font=\small}
]

\node[circle, draw=enspblue,  fill=enspblue!10,  minimum size=1.5cm, thick]
    (C) at ( 0, 5) {\textbf{Conf}};
\node[circle, draw=enspgreen, fill=enspgreen!10, minimum size=1.5cm, thick]
    (R) at (-4, 0) {\textbf{Rel}};
\node[circle, draw=navyblue,  fill=navyblue!10,  minimum size=1.5cm, thick]
    (O) at ( 4, 0) {\textbf{Opp}};

\draw[thick, enspblue!50]  (C) -- (R);
\draw[thick, enspgreen!50] (R) -- (O);
\draw[thick, navyblue!50]  (O) -- (C);

\node[draw=enspred, fill=enspred!8, rounded corners, thick,
      minimum width=3cm, minimum height=1.2cm] (bound) at (0, 2)
    {$\Conf \cdot \Rel \cdot \Opp \leq \dfrac{\Delta}{\lambda}$};

\draw[-{Stealth[length=2mm]}, thick, enspblue]  (C.south) -- (bound.north west);
\draw[-{Stealth[length=2mm]}, thick, enspgreen] (R.east)  -- (bound.south west);
\draw[-{Stealth[length=2mm]}, thick, navyblue]  (O.west)  -- (bound.south east);

\node[rotate=63,  font=\footnotesize, text=enspblue]  at (-2.3, 2.2)
    {Contrainte $\Delta$};
\node[rotate=-63, font=\footnotesize, text=navyblue]  at ( 2.3, 2.2)
    {Goulot $H(C)$};
\node[font=\footnotesize, text=enspgreen]              at (0,   -0.5)
    {Borne de Fano};

\node[above=0.3cm of C, font=\footnotesize, text=enspblue, align=center]
    {Protection\\de la vie priv\'ee};
\node[below left=0.3cm of R, font=\footnotesize, text=enspgreen, align=right]
    {Recouvrement\\du message};
\node[below right=0.3cm of O, font=\footnotesize, text=navyblue, align=left]
    {Interpr\'etation\\institutionnelle};

\end{tikzpicture}
\caption[Triangle du Trilemme CRO]{%
  Triangle du Trilemme CRO : trois propri\'et\'es de s\'ecurit\'e concurrentes.%
}\label{fig:trilemma_triangle}

\smallskip
\begin{minipage}{\linewidth}
  \small\textit{Note :} La borne centrale
  $\Conf \cdot \Rel \cdot \Opp \leq \Delta/\lambda$ limite l'optimisation simultan\'ee.
  Les \'etiquettes des ar\^etes indiquent le m\'ecanisme de contrainte principal
  reliant chaque paire de propri\'et\'es.
\end{minipage}

\end{figure}
\end{sloppypar}

% -----------------------------------------------------------------------
\subsection{Port\'ee du th\'eor\`eme et syst\`emes d\'eploy\'es}
\label{subsec:portee_trilemme}
% -----------------------------------------------------------------------
\begin{remark}[Port\'ee des \'etudes de cas]\label{rem:portee_cas_etude}
Sous l'Hypoth\`ese~\ref{ass:regime}, l'\'Etape~1 donne
$\Relrec \leq (1+\kappa/\lambda)/\lambda \approx 0{,}008$ \`a $\lambda = 128$.
Les syst\`emes d\'eploy\'es de r\'ef\'erence (Helios, zk-SNARKs, ECDSA),
dont la fiabilit\'e op\'erationnelle $\Relver$ est proche de~$1$, ne sont donc
\emph{pas} des instances du Th\'eor\`eme~\ref{thm:trilemme_cro_conditionnel},
car leurs m\'etriques CLO sont exprim\'ees sur l'\'echelle op\'erationnelle
de~\cite{MINKA2026CLO}, distincte des quantit\'es information-th\'eoriques des
D\'efinitions~\ref{def:confidentialite_clo}--\ref{def:opposabilite_clo}. Dans
la suite, ces syst\`emes servent \`a une \emph{calibration illustrative} des
ordres de grandeur des param\`etres institutionnels, non \`a une
validation de la borne. Le th\'eor\`eme s'applique aux protocoles \`a fuite
contr\^ol\'ee ($I(M;\Sigma|C) \leq \kappa/\lambda$), dont les engagements
parfaitement masquants et les preuves \`a divulgation nulle de connaissance
sont les exemples canoniques.
\end{remark}

% -----------------------------------------------------------------------
\subsection{R\'egimes op\'erationnels}
\label{subsec:regimes_operationnels}
% -----------------------------------------------------------------------
\begin{sloppypar}
Le Trilemme CRO classique distingue deux r\'egimes, s\'epar\'es par le seuil
$H^\sharp = (\Delta + 1 + \kappa/\lambda)/\alpha$ de la
Remarque~\ref{rem:activite_delta}. Ce seuil, d\'eriv\'e du th\'eor\`eme,
compare la richesse du contexte, pond\'er\'ee par la granularit\'e des
verdicts, \`a la capacit\'e d'extraction de l'institution.

\subsubsection{R\'egime \`a contexte rare~: $H(\mathcal{C}) \leq H^\sharp$}
\label{subsubsec:regime_faible_contexte}

La capacit\'e institutionnelle couvre la richesse du contexte~; le produit est
born\'e par l'in\'egalit\'e de Fano, $(1+\kappa/\lambda)/\lambda$. Avec les
param\`etres illustratifs de la calibration (Tableau~\ref{tab:validation_cro_classique}),
un contrat standardis\'e ($H(\mathcal{C}) \approx 20$ bits, $\alpha = 1$,
$\Delta \approx 35$ bits pour un juge g\'en\'eraliste~\cite{Mason07,Kerr04}, soit
$H^\sharp \approx 36$) et un scrutin de taille moyenne ($H(\mathcal{C}) \approx 40$
bits, $\alpha = 1{,}2$, $\Delta \approx 56$ bits pour un comit\'e de garants
sp\'ecialis\'es~\cite{Goodman15}, soit $H^\sharp \approx 48$) rel\`event de ce
r\'egime. La capacit\'e d'extraction n'y est pas la contrainte liante.

\subsubsection{R\'egime \`a contexte riche~: $H(\mathcal{C}) > H^\sharp$}
\label{subsubsec:regime_haute_contexte}
\label{subsubsec:regime_transition}

La granularit\'e des verdicts d\'epasse la capacit\'e d'extraction~; le produit
est born\'e par
\begin{equation}
\Conf \cdot \Rel \cdot \Opp \leq
\frac{1+\kappa/\lambda}{\lambda}\cdot\frac{\kappa/\lambda+\Delta}{\gamma},
\end{equation}
strictement plus petit que la borne de Fano. La conformit\'e financi\`ere via
zk-SNARKs~\cite{Groth16,BCTV14} en est l'exemple, avec $H(\mathcal{C}) \approx 70$
bits, $\alpha = 1{,}5$, $\Delta \approx 73$ bits pour un r\'egulateur
expert~\cite{BaselIII}, soit $H^\sharp \approx 49$ et $\gamma = 104$~; la borne
vaut environ $0{,}70$ fois la borne de Fano. C'est dans ce r\'egime que la
capacit\'e cognitive de l'institution, au-del\`a de la seule structure
cryptographique, plafonne le produit.

% -----------------------------------------------------------------------
\begin{table}[htbp]
\centering
\caption{Les deux r\'egimes op\'erationnels du Trilemme CRO classique}
\label{tab:regimes_operationnels}
\begin{tabular}{@{}lcc@{}}
\toprule
\textbf{Crit\`ere} & \textbf{Contexte rare} & \textbf{Contexte riche} \\
\midrule
Condition & $H(\mathcal{C}) \leq H^\sharp$ & $H(\mathcal{C}) > H^\sharp$ \\
Borne & $(1+\kappa/\lambda)/\lambda$ (Fano) & $\frac{1+\kappa/\lambda}{\lambda}\cdot\frac{\kappa/\lambda+\Delta}{\gamma}$ \\
Contrainte liante & fiabilit\'e de reconstruction & capacit\'e institutionnelle $\Delta$ \\
Construction de r\'ef\'erence & SOTP-Low / SOTP-Base & SOTP-High \\
Exemple (calibration) & ECDSA, Helios & zk-SNARKs \\
\bottomrule
\end{tabular}
\begin{tablenotes}
\item \footnotesize $H^\sharp = (\Delta + 1 + \kappa/\lambda)/\alpha$ (Remarque~\ref{rem:activite_delta}). Les exemples sont class\'es avec les param\`etres de calibration~; ils ne sont pas des instances du th\'eor\`eme (Remarque~\ref{rem:portee_cas_etude}).
\item \footnotesize Les constructions SOTP sont des constructions de r\'ef\'erence (chapitre~\ref{chap:constructions})~; leur optimalit\'e est une conjecture.
\end{tablenotes}
\end{table}
\end{sloppypar}

% -----------------------------------------------------------------------
\subsection{Calibration illustrative des param\`etres institutionnels}
\label{subsec:analyse_comparative}
% -----------------------------------------------------------------------
\begin{sloppypar}
Les syst\`emes d\'eploy\'es n'\'etant pas des instances du th\'eor\`eme
(Remarque~\ref{rem:portee_cas_etude}), leurs produits CLO ne peuvent ni
valider ni r\'efuter la borne. Ils permettent en revanche de calibrer, sur
l'\'echelle op\'erationnelle CLO, des ordres de grandeur plausibles de la
capacit\'e $\Delta$~\cite{Adida08,Groth16,Johnson01}.

% -----------------------------------------------------------------------
\begin{table}[htbp]
\centering
\caption{Calibration illustrative~: capacit\'es implicites sur l'\'echelle CLO}
\label{tab:validation_cro_classique}
\begin{tabular}{@{}lcccc@{}}
\toprule
\textbf{Syst\`eme} & $\boldsymbol{\Conf\cdot\Relver\cdot\Opp}$ & $\boldsymbol{\lambda}$ & $\boldsymbol{\Delta := \lambda \cdot \text{produit}}$ & $\boldsymbol{\Delta/\lambda}$ \\
\midrule
Helios~\cite{Adida08,Adida09UCL} & 0{,}437 & 128 & $\approx 56$ bits & 0{,}437 \\
zk-SNARKs~\cite{Groth16,BCTV14} & 0{,}572 & 128 & $\approx 73$ bits & 0{,}572 \\
ECDSA~\cite{Johnson01,NIST.FIPS.186-5} & 0{,}273 & 128 & $\approx 35$ bits & 0{,}273 \\
\bottomrule
\end{tabular}
\begin{tablenotes}
\item \footnotesize Les capacit\'es $\Delta$ sont \emph{d\'efinies} par $\Delta := \lambda \cdot (\Conf\cdot\Relver\cdot\Opp)$, c'est-\`a-dire par inversion de la forme publi\'ee $\Delta/\lambda$ sur l'\'echelle CLO~; les colonnes $\Delta/\lambda$ et produit co\"incident donc par construction.
\item \footnotesize Ces syst\`emes ne sont pas des instances du Th\'eor\`eme~\ref{thm:trilemme_cro_conditionnel} (Remarque~\ref{rem:portee_cas_etude}). Voir la note m\'ethodologique de la sous-section~\ref{subsec:note_methodo}.
\end{tablenotes}
\end{table}
% -----------------------------------------------------------------------

Chaque capacit\'e implicite peut \^etre rapproch\'ee d'une d\'ecomposition en
ressources cognitives et organisationnelles, d\'etaill\'ee au
chapitre~\ref{chap:experimentation}
(Tableau~\ref{tab:decomposition_delta})~:

\begin{itemize}
\item \textbf{Helios}~\cite{Adida08,Goodman15}, avec $\Delta \approx 56$ bits pour un comit\'e
\'electoral sp\'ecialis\'e (expertise, pr\'ec\'edents, proc\'edures, collaboration)~\cite{Ericsson95,Simon82}.

\item \textbf{zk-SNARKs}~\cite{Groth16,BCTV14}, avec $\Delta \approx 73$ bits pour une institution
financi\`ere experte (cadres r\'eglementaires, analyse de sch\'emas, \'equipes
forensiques, collaboration multi-agences)~\cite{BaselIII}.

\item \textbf{ECDSA}~\cite{Johnson01,Mason07}, avec $\Delta \approx 35$ bits pour un juge de
juridiction g\'en\'erale (pr\'ec\'edents, proc\'edures, expertise g\'en\'erale)~\cite{Kerr04,Casey11}.
\end{itemize}

Les capacit\'es implicites ainsi calibr\'ees ($35$--$73$ bits) tombent dans
l'intervalle de quelques dizaines de bits sugg\'er\'e par la litt\'erature sur la
rationalit\'e born\'ee~\cite{Miller56,Simon82,Ericsson95,Cowan2001}. Il s'agit
d'une v\'erification de plausibilit\'e des ordres de grandeur, distincte d'une validation du th\'eor\`eme. Cette calibration est \textbf{r\'etrospective et
illustrative}~; elle ne constitue pas une preuve empirique des hypoth\`eses
institutionnelles~\cite{Tsybakov2009,Kahneman11}.
\end{sloppypar}

% -----------------------------------------------------------------------
\subsection{Positionnement par rapport au r\'egime quantique}
\label{subsec:transition_postquantique}
% -----------------------------------------------------------------------
\begin{sloppypar}
Le Trilemme CRO classique \'etabli dans ce chapitre s'applique aux adversaires
computationnels classiques (PPT)~\cite{Goldreich01,KatzLindell14}. La migration vers les standards post-quantiques
NIST soul\`eve une question critique, celle de savoir si ces bornes tiennent face \`a des
adversaires quantiques~\cite{NISTPQC,Watrous2009,Boneh2011,Zhandry}.

La section~\ref{sec:cro_pq} \'etendra ce cadre en introduisant~:

\begin{itemize}
\item La capacit\'e d'extraction quantique $\Delta_q = \Delta + \delta_{\text{enh}} - \delta_{\text{degr}}$
combinaison de la capacit\'e classique et des modifications induites par les
techniques quantiques~\cite{NielsenChuang2010,Wilde2017}.

\item Le coefficient de perte d'interpr\'etabilit\'e quantique $\eta_q$ quantifiant
la d\'egradation s\'emantique dans les canaux quantiques (d\'ecoh\'erence, erreurs de portes)~\cite{Zurek2003,Wootters1982}.

\item La borne post-quantique~:
\begin{equation}
\Conf^q \cdot \Rel^q \cdot \Opp^q \leq \frac{\Delta_q + \eta_q \cdot H(C)}{\lambda} + O(\lambda^{-2})
\end{equation}
\end{itemize}

La borne explicite restera port\'ee par la capacit\'e quantique $\Delta_q$ et par
$\gamma_q$~; le terme $\eta_q \cdot H(C)/\lambda$ n'appara\^it que dans une
majoration affaiblie, commode pour les ordres de grandeur. Le bruit quantique
agit comme une r\'eduction de l'information interpr\'etable, dont une
formalisation par contraction est propos\'ee en perspective
(chapitre~\ref{chap:resultats_discussions}).
\end{sloppypar}

% -----------------------------------------------------------------------
\begin{boiteArbitre}
L'arbitre referme le dossier du Trilemme CRO classique et termine ainsi l'\'etape \textbf{B}. Cette section lui a
appris un r\'esultat contre-intuitif, puisque la limite se trouve en lui plut\^ot
que dans la cryptographie d\`es que ses verdicts sont fins.

Quatre contributions structurent ce qu'il retient. Le mod\`ele \'etendu
d'abord, avec trois param\`etres mesurables ($\Delta$, $\alpha$, $\beta$) qui
traduisent ses contraintes cognitives en grandeurs quantifiables. Les
hypoth\`eses ensuite, explicites et honn\^etes sur ce qu'elles ne prouvent pas
encore. Le th\'eor\`eme lui-m\^eme, $\Conf \cdot \Rel \cdot \Opp \leq
\Delta/(\gamma\lambda) + O(\lambda^{-2})$, puis sa forme raffin\'ee, qui lui dit
exactement quand sa capacit\'e d'extraction devient la contrainte liante. La
distinction des deux r\'egimes enfin, qui transforme ce r\'esultat n\'egatif en
guide de conception.

Il est surtout marqu\'e par la franchise sur la port\'ee. Helios, les
zk-SNARKs et ECDSA ne sont pas des instances du th\'eor\`eme et calibrent
seulement les ordres de grandeur~; les capacit\'es $\Delta$ qui en
d\'ecoulent ($35$ \`a $73$ bits) restent plausibles au regard de la
litt\'erature cognitive.

La borne est n\'ecessaire sans \^etre encore prouv\'ee serr\'ee de mani\`ere
g\'en\'erale. La section suivante l'\'etendra au r\'egime quantique.

Il veut voir si les constructions tiennent ce que le th\'eor\`eme promet.
\end{boiteArbitre}